\sgaStableTarget{sga6.e14.eq.6-2}@@M0@@
where @@M1@@ is the virtual relative dimension of @@M2@@. Assuming this homomorphism defined, and taking account of the Riemann--Roch formula in the form of Exp. VIII with values in @@M3@@, one finds that the validity of the @@M4@@-adic Riemann--Roch formula for the morphism @@M5@@ is exactly equivalent to the commutativity of the square
\sgaStableTarget{sga6.e14.eq.6-3}@@M6@@
where the horizontal arrows are defined from (6.1) by passing to associated graded groups and dividing so as to take Chern classes to Chern classes, and the vertical arrows are respectively the arrow of Exp. VIII and the Gysin homomorphism in cohomology (assumed given).

\sgaStableTarget{sga6.e14.para.6-2}\textbf{6.2.} Observe that the Gysin homomorphism is in any case defined when @@M7@@ is smooth over a field @@M8@@ (SGA 5 IV), using the relative cohomological purity theorem over @@M9@@; the same construction applies when @@M10@@ is excellent, regular and of characteristic zero, by the absolute cohomological purity theorem of SGA 4 XIX, and it will apply automatically to every excellent regular @@M11@@, without restriction on the characteristic, once the absolute cohomological purity theorem has been proved in full generality. Finally, the same construction applies, by the relative cohomological purity theorem (SGA 4 XVI 3.7), when @@M12@@ and @@M13@@ are both smooth over the same scheme @@M14@@ and @@M15@@ is an @@M16@@-morphism. One should conjecture that, in each of these cases in which the Gysin homomorphism is defined, diagram (6.3) is indeed commutative. The only case in which this assertion is currently verified is that in which @@M17@@ and @@M18@@ are both smooth over a field @@M19@@; it then follows from the commutativity of the diagram analogous to (6.3):
\sgaStableTarget{sga6.e14.eq.6-4}@@M20@@
